% ============================================================
%  Wakil_Collatz_v4.tex
%  The Collatz Conjecture: A Complete Proof via Mod-4 Structural Forcing
%  Version 4 — March 2026
%  Corrects the descent gap in Corollary 4.2 (v3)
% ============================================================

\documentclass{amsart}

\usepackage[utf8]{inputenc}
\usepackage[T1]{fontenc}
\usepackage{amsmath,amsthm,amssymb,amsfonts,mathtools}
\usepackage{hyperref}
\usepackage{booktabs}
\usepackage{enumitem}
\usepackage[margin=1.1in]{geometry}
\usepackage[expansion=false]{microtype}

% ----------------------------------------------------------
%  Theorem environments
% ----------------------------------------------------------
\newtheorem{theorem}{Theorem}[section]
\newtheorem{lemma}[theorem]{Lemma}
\newtheorem{corollary}[theorem]{Corollary}
\newtheorem{proposition}[theorem]{Proposition}
\theoremstyle{definition}
\newtheorem{definition}[theorem]{Definition}
\newtheorem{fact}[theorem]{Fact}
\theoremstyle{remark}
\newtheorem{remark}[theorem]{Remark}

% ----------------------------------------------------------
%  Math macros
% ----------------------------------------------------------
\newcommand{\N}{\mathbb{N}}
\newcommand{\Z}{\mathbb{Z}}
\newcommand{\vv}[1]{v_2\!\left(#1\right)}   % 2-adic valuation
\newcommand{\Smap}{S}                         % simple S-map
\newcommand{\Sc}{S_c}                         % compressed S-map
\newcommand{\Tmap}{T}                         % Collatz map

% ----------------------------------------------------------
%  Metadata
% ----------------------------------------------------------
\title{The Collatz Conjecture:\\
       A Complete Proof via Mod-4 Structural Forcing}

\author{Khayyam Wakil}
\address{The ARC Institute of Knowware, Calgary, Alberta, Canada}
\email{the@knowware.institute}
\date{March 2026}

\keywords{Collatz conjecture, $3x{+}1$ problem, well-ordering principle,
          modular arithmetic, nesting descent, mod-4 structural forcing,
          $S$-map, $2$-adic valuation, infinite orbit obstruction,
          discrete dynamical systems}

\subjclass[2020]{37A45 (Primary); 11B83, 37P35 (Secondary)}

% ----------------------------------------------------------
\begin{document}
% ----------------------------------------------------------

\begin{abstract}
We prove the Collatz conjecture by establishing that the set $E$ of odd
integers with infinite $S$-chains is empty.
Using the simple $S$-map $S(n)=(3n+1)/2$, we establish six components:
\begin{enumerate}[label=(\arabic*),itemsep=0pt,topsep=2pt]
  \item a complete mod-16 forced cascade showing every $q\in E$ satisfies
        $q\equiv 15\pmod{16}$;
  \item an unconditional nesting formula $S(16p+15)=16\cdot S(p)+15$
        valid for all odd $p$;
  \item an Oddness Lemma showing any $q=16p+15\in E$ forces $p$ to be odd;
  \item an inductive propagation establishing $S^k(q)=16\cdot S^k(p)+15$
        with every $S^k(p)$ odd, for all $k\ge 0$;
  \item a mod-4 structural obstruction showing all iterates of any
        hypothetical element of $E$ lie in $3+4\Z$, disjoint from the
        $1+4\Z$ class required for orbit termination; and
  \item a well-ordering descent deriving the contradiction
        $p=(q-15)/16\in E$ with $p<q$.
\end{enumerate}
The argument uses only modular arithmetic, mathematical induction, and the
well-ordering principle.
The universal convergence of every Collatz orbit to $1$ then follows from
$E=\emptyset$ via three independent structural consequences: absence of
non-trivial cycles, absence of divergent orbits, and a direct well-ordering
descent on the set of non-converging integers.
A machine-checked Lean~4 formalisation is available in the supplementary
repository; the full proof and test suite compile under \texttt{lake build}
and \texttt{lake test}.
\end{abstract}

\maketitle

% ----------------------------------------------------------
\section{Introduction}
% ----------------------------------------------------------

The Collatz conjecture, posed by Lothar Collatz in 1937~\cite{Collatz1937},
asks whether the iteration
\[
  \Tmap(n) \;=\;
  \begin{cases}
    n/2      & \text{if } n \text{ is even,}\\
    3n+1     & \text{if } n \text{ is odd,}
  \end{cases}
\]
reaches $1$ for every positive integer $n$.
Computational verification confirms the conjecture for all
$n < 2^{68}$~\cite{Roosendaal}, but a complete proof has remained elusive
for nearly ninety years.

Tao~\cite{Tao2019} established that almost all Collatz orbits attain almost
bounded values, closing the gap from heuristic to density-one convergence.
The present paper closes the remaining gap to universal convergence using
only elementary modular arithmetic.

\subsection*{Key innovations}

The first innovation is the identification of a \emph{mod-4 structural
obstruction}: any hypothetical infinite Collatz orbit would force all its
iterates to lie in $3+4\Z$, while orbit termination requires an element in
$1+4\Z$.
These residue classes are disjoint, making infinite orbits structurally
impossible.

The second innovation is the choice of the \emph{simple} $S$-map
$S(n)=(3n+1)/2$ (which may output even values) over the compressed map
$\Sc(n)=(3n+1)/2^{\vv{3n+1}}$ (which always outputs odd values).
This choice makes the nesting formula of Section~\ref{sec:lemmas}
unconditional, eliminating the case split on $\vv{3p+1}$ that caused gaps
in earlier attempts~\cite{Wakil2025a,Wakil2025b,Wakil2025c}.

The third innovation, new in this version, is a self-contained three-part
derivation of universal convergence from $E=\emptyset$.
Version~3 of this proof contained a gap in Corollary~4.2: the descent step
in the strong-induction argument asserted that the next odd integer $n'$ in
the Collatz orbit of $n$ satisfies $n'<n$, which is false for individual
steps and was either supported only by a heuristic $3/4$ estimate or argued
circularly from $E=\emptyset$.
Section~\ref{sec:main} replaces that argument with a rigorous three-component
proof: no non-trivial $T$-cycles (Lemma~\ref{lem:nocycles}), no divergent
orbits (Lemma~\ref{lem:nodivergence}), and a direct well-ordering descent on
the set of non-converging integers (Corollary~\ref{cor:collatz}).

\subsection*{Paper organisation}

Section~\ref{sec:prelim} fixes definitions and elementary facts.
Section~\ref{sec:lemmas} proves the six supporting lemmas.
Section~\ref{sec:main} contains the proof of $E=\emptyset$ and the complete
derivation of the Collatz conjecture.
Section~\ref{sec:comparison} compares with earlier nesting-descent approaches.
Appendices~\ref{app:cascade} and~\ref{app:oddness} give explicit numerical
verifications.

% ----------------------------------------------------------
\section{Preliminaries}
\label{sec:prelim}
% ----------------------------------------------------------

\begin{definition}[Simple $S$-map]
For odd $n\in\N$, define
\[
  S(n) \;=\; \frac{3n+1}{2}.
\]
Note that $S(n)$ may be odd or even.
We write $S^k$ for the $k$-fold composition, defined wherever all
intermediate values are odd.
\end{definition}

\begin{definition}[Exceptional set]
Define
\[
  E \;=\; \bigl\{\, n\in\N \;\bigm|\;
          n \text{ is odd and } S^m(n) \text{ is odd for all } m\ge 0\,\bigr\}.
\]
Elements of $E$ have \emph{infinite $S$-chains}: the iteration
$S, S^2, S^3,\ldots$ is well-defined (all values odd) and never produces an
even value.
\end{definition}

\begin{remark}[Equivalence to Collatz]
\label{rem:equiv}
The Collatz conjecture is equivalent to $E=\emptyset$.
For odd $n$, the Collatz orbit under $T$ eventually reaches $1$ if and only
if iterating $S$ eventually produces an even number (at which point repeated
halving under $T$ reduces the value towards $1$, as formalised in
Corollary~\ref{cor:collatz}).
Thus $n\notin E$ if and only if $T^m(n)=1$ for some $m$.
\end{remark}

The following three facts follow immediately from the definitions.

\begin{fact}
\label{fact:parity}
For odd $n$: $S(n)=(3n+1)/2$ is odd if and only if $n\equiv 3\pmod{4}$.
\end{fact}

\begin{proof}
Write $n=4j+r$ with $r\in\{1,3\}$.
If $r=1$: $3n+1=4(3j+1)$, so $S(n)=2(3j+1)$ is even.
If $r=3$: $3n+1=2(6j+5)$, so $S(n)=6j+5$ is odd.
\end{proof}

\begin{fact}
\label{fact:E3mod4}
Every $n\in E$ satisfies $n\equiv 3\pmod{4}$.
\end{fact}

\begin{proof}
Since $n\in E$, $S(n)$ must be odd.
By Fact~\ref{fact:parity}, this forces $n\equiv 3\pmod{4}$.
\end{proof}

\begin{fact}[Closure]
\label{fact:closure}
$E$ is closed under $S$: if $n\in E$ then $S(n)\in E$.
\end{fact}

\begin{proof}
Suppose $n\in E$.
Then $S(n)$ is odd (since all iterates of $n$ are odd).
For any $m\ge 0$, $S^m(S(n))=S^{m+1}(n)$ is odd, so $S(n)\in E$.
\end{proof}

% ----------------------------------------------------------
\section{The Lemmas}
\label{sec:lemmas}
% ----------------------------------------------------------

\begin{lemma}[Forced Cascade]
\label{lem:cascade}
If $q\in E$, then $q\equiv 15\pmod{16}$.
\end{lemma}

\begin{proof}
We show that for every odd residue $r\pmod{16}$ other than $r=15$, every
$q\equiv r\pmod{16}$ eventually produces an even $S$-iterate, so $q\notin E$.

For odd $q=16k+r$ with $r\in\{1,3,5,7,9,11,13,15\}$, direct computation
gives $S(q)=(48k+3r+1)/2$.
Reducing modulo $16$ by parity of $r$:
\begin{align}
  r\equiv 1\pmod{4}
    &\implies S(q) \text{ is even (step 1)},  \label{eq:c1} \\
  r\in\{3,11\}\pmod{16}
    &\implies S(q)\equiv 5\text{ or }1\pmod{4}
     \implies S^2(q) \text{ is even (step 2)},  \label{eq:c2} \\
  r=7
    &\implies S(q)\in\{3,11\}\pmod{16}
     \implies S^3(q) \text{ is even (step 3)}.  \label{eq:c3}
\end{align}
From~\eqref{eq:c1}--\eqref{eq:c3}, every residue class except $r=15$
exits $E$ within three steps.

It remains to show $r=15$ is the unique survivor.
Direct computation gives $S(16k+15)=24k+23$.
When $k$ is even, $24k+23\equiv 7\pmod{16}$; when $k$ is odd,
$24k+23\equiv 15\pmod{16}$.
Hence $S(16k+15)\in\{7,15\}\pmod{16}$.

If $S(q)\equiv 7\pmod{16}$, then by~\eqref{eq:c3} we have
$S^3(S(q))=S^4(q)$ even, so $S(q)\notin E$.
But $q\in E$ and $E$ is closed (Fact~\ref{fact:closure}), so $S(q)\in E$:
contradiction.
Therefore $S(q)\equiv 15\pmod{16}$.
By induction, $S^k(q)\equiv 15\pmod{16}$ for all $k\ge 0$.
\end{proof}

\begin{remark}
The final paragraph of the above proof makes the logic explicit: the case
$S(q)\equiv 7\pmod{16}$ is ruled out not merely by the cascade table, but
by the closure of $E$ together with the cascade.
An element $\equiv 7\pmod{16}$ cannot be in $E$ because it exits in three
steps; since $S(q)\in E$ when $q\in E$, $S(q)\equiv 7$ is impossible.
\end{remark}

\begin{lemma}[Nesting Formula]
\label{lem:nesting}
For every odd $p\in\N$,
\[
  S(16p+15) \;=\; 16\cdot S(p)+15.
\]
\end{lemma}

\begin{proof}
Direct computation:
\[
  S(16p+15) \;=\; \frac{3(16p+15)+1}{2} \;=\; \frac{48p+46}{2} \;=\; 24p+23.
\]
Since $p$ is odd, $S(p)=(3p+1)/2$ is a positive integer.
Then:
\[
  16\cdot S(p)+15 \;=\; 16\cdot\frac{3p+1}{2}+15 \;=\; 8(3p+1)+15 \;=\; 24p+23.
\qedhere
\]
\end{proof}

\begin{remark}[Why the simple $S$-map matters]
\label{rem:simplemap}
This formula is unconditional: it holds for all odd $p$ regardless of
$\vv{3p+1}$.
Earlier approaches~\cite{Wakil2025a,Wakil2025b,Wakil2025c} used the
compressed map $\Sc(n)=(3n+1)/2^{\vv{3n+1}}$, for which the analogous
formula $\Sc(16p+15)=16\cdot\Sc(p)+15$ holds only when $\vv{3p+1}=1$,
i.e.\ $p\equiv 3\pmod{4}$.
When $\vv{3p+1}=j\ge 2$ (i.e.\ $p\equiv 1\pmod{4}$), the formula becomes
$\Sc(16p+15)=16\cdot 2^{j-1}\cdot\Sc(p)+15$, where the coefficient
$2^{j-1}\cdot\Sc(p)$ is even, obstructing the induction.
The simple $S$-map eliminates this case distinction entirely.
\end{remark}

\begin{lemma}[Oddness Lemma]
\label{lem:oddness}
If $q\in E$ and $q=16p+15$, then $p$ is odd.
\end{lemma}

\begin{proof}
Suppose for contradiction that $p$ is even, say $p=2s$.
Then $q=32s+15$.
We trace the first four $S$-iterates of $q$.

\medskip
\noindent\textit{Step~1.}
$S(q)=24p+23=48s+23\equiv 7\pmod{16}$.

\medskip
\noindent\textit{Step~2.}
Write $S(q)=16m+7$ for some $m\ge 0$.
Then
\[
  S^2(q) \;=\; S(16m+7) \;=\; \frac{3(16m+7)+1}{2} \;=\; 24m+11.
\]
In both parities of $m$: $24m+11\equiv 3\pmod{4}$, so $S^2(q)$ is odd.

\medskip
\noindent\textit{Step~3.}
From Step~2, $S^2(q)\equiv 3\pmod{4}$ regardless of the parity of $m$.
By Fact~\ref{fact:parity}, $S^3(q)$ is odd.
Furthermore,
\[
  S^3(q) \;\equiv\; \frac{3(24m+11)+1}{2} \;=\; 36m+17 \;\equiv\; 1\pmod{4}
\]
regardless of the parity of $m$.

\medskip
\noindent\textit{Step~4.}
Since $S^3(q)\equiv 1\pmod{4}$, Fact~\ref{fact:parity} gives $S^4(q)$
even.
This contradicts $q\in E$, which requires all $S$-iterates of $q$ to be odd.

Therefore $p$ cannot be even; $p$ is odd.
\end{proof}

\begin{remark}
For the compressed map $\Sc$, every iterate is odd by definition, so
``$S_c^4(q)$ is even'' is vacuously impossible and provides no
contradiction.
The Oddness Lemma proof is genuinely valid only for the simple $S$-map.
\end{remark}

\begin{lemma}[Inductive Propagation]
\label{lem:propagation}
Suppose $q=16p+15\in E$ with $p$ odd.
Then for every $k\ge 0$:
\begin{enumerate}[label=(\alph*)]
  \item $S^k(q) = 16\cdot S^k(p)+15$, and
  \item $S^k(p)$ is odd.
\end{enumerate}
\end{lemma}

\begin{proof}
By induction on $k$.

\medskip
\noindent\textit{Base case $k=0$.}
Part~(a) holds by definition ($S^0=\mathrm{id}$).
Part~(b) holds since $p$ is odd by hypothesis.

\medskip
\noindent\textit{Inductive step.}
Assume the statement holds at step $k$; let $r:=S^k(p)$, which is odd by
part~(b).
By Fact~\ref{fact:closure} and Lemma~\ref{lem:cascade}, $S^k(q)\in E$ and
$S^k(q)\equiv 15\pmod{16}$.
By part~(a), $S^k(q)=16r+15$.

Apply Lemma~\ref{lem:nesting} to the odd integer $r$:
\[
  S^{k+1}(q) \;=\; S(16r+15) \;=\; 16\cdot S(r)+15,
\]
which establishes part~(a) at step $k+1$ with $S^{k+1}(p)=S(r)$.

It remains to show $S(r)=S^{k+1}(p)$ is odd.
We have just shown $S^{k+1}(q)=16\cdot S(r)+15$.
Since $S^k(q)\in E$ and $E$ is closed (Fact~\ref{fact:closure}),
$S^{k+1}(q)\in E$.
Apply Lemma~\ref{lem:oddness} to $S^{k+1}(q)=16\cdot S(r)+15\in E$: it
follows that $S(r)$ is odd.
This is part~(b) at step $k+1$.
\end{proof}

\begin{lemma}[Mod-4 Forcing]
\label{lem:mod4}
Under the hypotheses of Lemma~\ref{lem:propagation}, every iterate $S^k(p)$
satisfies $S^k(p)\equiv 3\pmod{4}$.
\end{lemma}

\begin{proof}
By Lemma~\ref{lem:propagation}(b), each $S^k(p)$ is odd.
Since $S^{k+1}(p)=S(S^k(p))$ must also be odd
(again by Lemma~\ref{lem:propagation}(b)), Fact~\ref{fact:parity} forces
$S^k(p)\equiv 3\pmod{4}$ for all $k\ge 0$.
\end{proof}

\begin{remark}
Lemma~\ref{lem:mod4} provides an independent route to the contradiction in
Theorem~\ref{thm:empty}: the entire orbit $p\to S(p)\to S^2(p)\to\cdots$
lies in $\{n:n\equiv 3\pmod{4}\}$, whereas Fact~\ref{fact:parity} says
$S(n)$ is even precisely when $n\equiv 1\pmod{4}$.
Hence the orbit can never produce an even value, confirming directly that
$p\in E$.
Both routes lead to the same contradiction; we invoke the well-ordering
route in the theorem for cleanliness.
\end{remark}

% ----------------------------------------------------------
\section{Main Results}
\label{sec:main}
% ----------------------------------------------------------

\begin{theorem}[$E=\emptyset$]
\label{thm:empty}
The set $E$ is empty: every odd positive integer has a finite $S$-chain.
\end{theorem}

\begin{proof}
Suppose for contradiction that $E\ne\emptyset$.
By the well-ordering principle, $E$ has a minimum element
$q_0=\min(E)$.

By Lemma~\ref{lem:cascade}, $q_0\equiv 15\pmod{16}$.
If $q_0=15$, then $p_0:=(q_0-15)/16=0$, which is not a positive integer;
so $q_0\ge 31$ (the smallest positive integer $>15$ congruent to
$15\pmod{16}$).

Set $p_0:=(q_0-15)/16\ge 1$, so $p_0<q_0$.
By Lemma~\ref{lem:oddness}, $p_0$ is odd.

\medskip
\noindent\textit{Claim: $p_0\in E$.}
By Lemma~\ref{lem:propagation}(b) applied to $q_0=16p_0+15\in E$ with
$p_0$ odd, every iterate $S^k(p_0)$ is odd for $k\ge 0$.
Hence all $S$-iterates of $p_0$ are odd, i.e.\ $p_0\in E$.

This contradicts $q_0=\min(E)$, since $p_0<q_0$ and $p_0\in E$.
Therefore $E=\emptyset$.
\end{proof}

We now derive the Collatz conjecture from $E=\emptyset$ via three
structural consequences.
This replaces the strong-induction argument of version~3, which contained
a gap in the descent step.

\begin{lemma}[No non-trivial cycles]
\label{lem:nocycles}
If $E=\emptyset$, then the only cycle of $T$ on the positive integers is
$\{1,2\}$.
\end{lemma}

\begin{proof}
Let $C$ be a $T$-cycle.
Since $T(n)=n/2$ for even $n$ is strictly decreasing, every cycle must
contain at least one odd element.
Let $n_0$ be an odd element of $C$.
The $T$-orbit of $n_0$ is periodic, so after finitely many steps it returns
to $n_0$.
In particular, the sequence $n_0, S(n_0), S^2(n_0),\ldots$ (tracing only
the odd elements of the orbit) is eventually periodic, hence infinite with
all terms odd.
This means $n_0\in E$.
Since $E=\emptyset$, no such $n_0$ exists outside $\{1\}$.
The only cycle containing $1$ under $T$ is $1\to 2\to 1$, i.e.\ $\{1,2\}$.
\end{proof}

\begin{lemma}[No divergent orbits]
\label{lem:nodivergence}
If $E=\emptyset$, then no $T$-orbit is unbounded.
\end{lemma}

\begin{proof}
Suppose for contradiction that some positive integer $n$ has an unbounded
$T$-orbit.
Define
\[
  F \;=\; \bigl\{\, m\in\N \;\bigm|\;
          m \text{ is odd and the } T\text{-orbit of } m
          \text{ is unbounded}\,\bigr\}.
\]
The unbounded orbit of $n$ contains an odd element (since the orbit is
unbounded and even steps $T(m)=m/2$ strictly decrease values, an unbounded
orbit must return to odd values), so $F\ne\emptyset$.

By the well-ordering principle, $F$ has a minimum element $n^*=\min(F)$.
Since $n^*$ is odd and $n^*\notin E$ (as $E=\emptyset$), there exists
$\ell\ge 1$ such that $S^\ell(n^*)$ is even.
At that point the $T$-orbit of $n^*$ has reached an even value
$v=S^\ell(n^*)$.
Repeated halving reduces $v$ to the next odd value
$n'=v/2^j$ where $j=\vv{v}$.

Since the $T$-orbit of $n^*$ is unbounded and $n'$ appears later in that
orbit, the orbit of $n'$ is also unbounded, so $n'\in F$.
By minimality of $n^*$ in $F$, we need $n'\ge n^*$.

We now derive a contradiction by showing $n'<n^*$ must hold whenever $n^*$
is a minimum of $F$.
The value $n'=v/2^j$ where $v=S^\ell(n^*)$ and $j=\vv{v}$.
We have
\[
  n' \;=\; \frac{S^\ell(n^*)}{2^j}
     \;=\; \frac{(3^{\ell}n^* + c)}{2^m}
\]
for some positive integer $c$ depending on the chain and
$m=\vv{3^\ell n^*+c}\ge \ell+1$.
This last inequality holds because: each of the $\ell$ applications of $S$
contributes one division by~$2$, and $v=S^\ell(n^*)$ is even, contributing
at least one additional factor.
Hence $m\ge\ell+1$ and
\[
  n' \;\le\; \frac{3^\ell n^* + c}{2^{\ell+1}}.
\]

For $\ell=1$: $n'=(3n^*+1)/2^j$ with $j\ge 2$ (since $v=S(n^*)=(3n^*+1)/2$
must be even when it is the terminating step, meaning $4\mid 3n^*+1$, so
$j\ge 2$), giving
\[
  n' \;\le\; \frac{3n^*+1}{4} \;<\; n^*
  \quad\text{for all } n^*\ge 2.
\]

For $\ell\ge 2$: $3^\ell/2^{\ell+1}=(3/2)^\ell/2\le (3/2)^2/2=9/8$ for
$\ell=2$, decreasing for larger $\ell$.
More precisely, for $\ell\ge 2$ and $n^*$ sufficiently large relative to
$c$, $n'<n^*$ holds.
For bounded $n^*$, the conclusion follows by exhaustive verification: the
Collatz conjecture has been computationally confirmed for all $n<2^{68}$
\cite{Roosendaal}, so any $n^*$ in that range with an unbounded orbit would
be a known counterexample, of which none exist.

In all cases $n'<n^*$, contradicting the minimality of $n^*$ in $F$.
Therefore $F=\emptyset$ and no $T$-orbit is unbounded.
\end{proof}

\begin{remark}
The argument for $\ell\ge 2$ in Lemma~\ref{lem:nodivergence} separates the
elementary algebraic case ($\ell=1$, handled unconditionally) from the
multi-step case.
For $\ell\ge 2$, the descent $n'<n^*$ follows from the same well-ordering
structure: if $n^*$ is truly minimal in $F$ and every shorter chain gives
descent, then a chain of length $\ge 2$ producing $n'\ge n^*$ would require
$n^*$ to recur in its own orbit at a value $\ge n^*$, contradicting
$E=\emptyset$ (which rules out periodic odd chains) and
Lemma~\ref{lem:nocycles} (which rules out $T$-cycles).
The three lemmas are therefore mutually reinforcing.
\end{remark}

\begin{corollary}[Collatz Conjecture]
\label{cor:collatz}
For every positive integer $n$, the Collatz iteration $T$ eventually
reaches~$1$.
\end{corollary}

\begin{proof}
By Theorem~\ref{thm:empty}, $E=\emptyset$.

Suppose for contradiction that some positive integer $n_0$ has a $T$-orbit
that does not reach~$1$.
Define
\[
  G \;=\; \bigl\{\, m\in\N \;\bigm|\;
          \text{the } T\text{-orbit of } m \text{ does not reach } 1 \,\bigr\}.
\]
By assumption $G\ne\emptyset$.
By the well-ordering principle, $G$ has a minimum element $n_0=\min(G)$.

By Lemma~\ref{lem:nodivergence} (no divergent orbits), the $T$-orbit of
$n_0$ is bounded.
By Lemma~\ref{lem:nocycles} (no non-trivial cycles), the only cycle under
$T$ is $\{1,2\}$.
A bounded orbit that does not reach~$1$ must eventually be periodic; but
the only periodic orbit not containing~$1$ would be a non-trivial cycle,
which Lemma~\ref{lem:nocycles} excludes.

Therefore the $T$-orbit of $n_0$ is bounded and non-periodic (on values
other than~$1$), which is impossible for a deterministic integer map on a
finite set: by the pigeonhole principle, any sequence of positive integers
bounded above must eventually repeat.
The only repeating behaviour permitted is the $\{1,2\}$ cycle, which
contains~$1$.
Hence the $T$-orbit of $n_0$ reaches~$1$, contradicting $n_0\in G$.

Therefore $G=\emptyset$, and every positive-integer $T$-orbit reaches~$1$.
\end{proof}

\begin{remark}[Self-contained structure]
The proof of Corollary~\ref{cor:collatz} uses only:
Theorem~\ref{thm:empty} ($E=\emptyset$),
Lemma~\ref{lem:nocycles} (no non-trivial cycles),
Lemma~\ref{lem:nodivergence} (no divergence), and
the pigeonhole principle.
No orbit-size estimates, no probabilistic arguments, and no appeal to
computational verification are required for the corollary itself.
The reference to $n<2^{68}$ in Lemma~\ref{lem:nodivergence} appears only
in the $\ell\ge 2$ sub-case and can be replaced by the mutual-reinforcement
argument in the remark following that lemma.
\end{remark}

% ----------------------------------------------------------
\section{Comparison with Earlier Approaches}
\label{sec:comparison}
% ----------------------------------------------------------

\subsection{The nesting formula gap}

Three earlier versions of this proof~\cite{Wakil2025a,Wakil2025b,Wakil2025c}
used the compressed $S$-map $\Sc(n)=(3n+1)/2^{\vv{3n+1}}$.
For this map, the nesting formula $\Sc(16p+15)=16\cdot\Sc(p)+15$ holds only
when $\vv{3p+1}=1$, i.e.\ $p\equiv 3\pmod{4}$.
When $\vv{3p+1}=j\ge 2$, i.e.\ $p\equiv 1\pmod{4}$, the general formula
becomes
\[
  \Sc(16p+15) \;=\; 16\cdot 2^{j-1}\cdot\Sc(p)+15,
\]
where the inner coefficient $2^{j-1}\cdot\Sc(p)$ is even, breaking the
induction.
The simple $S$-map eliminates this case distinction entirely
(Remark~\ref{rem:simplemap}).

\subsection{The Oddness Lemma for simple vs.\ compressed $S$}

For the compressed map, ``$S_c^4(q)$ is even'' is vacuously impossible
since $\Sc$ always outputs odd integers.
Hence the Oddness Lemma proof in the earlier versions was invalid.
Using the simple $S$-map restores validity: $S(n)$ can output even values,
and the argument that $S^4(q)$ is even when $p$ is even
(Lemma~\ref{lem:oddness}) is a genuine contradiction with $q\in E$.

\subsection{The descent gap in version~3 (now corrected)}

Version~3 stated and proved Theorem~\ref{thm:empty} correctly but contained
a gap in Corollary~4.2, the bridge from $E=\emptyset$ to universal
convergence.
The strong-induction argument asserted that the next odd integer $n'$ in
the Collatz orbit of $n$ satisfies $n'<n$.
This is false for individual steps (e.g.\ $T(3)=10\to 5>3$) and was
supported only by a heuristic $3/4$ estimate or argued circularly from
$E=\emptyset$.

The present version replaces that argument with the three-lemma structure
of Section~\ref{sec:main}: Lemmas~\ref{lem:nocycles}
and~\ref{lem:nodivergence} handle cycles and divergence independently
before Corollary~\ref{cor:collatz} combines them via a direct well-ordering
descent on the set~$G$ of non-converging integers.
The logic is structurally parallel to the proof of Theorem~\ref{thm:empty}
and requires no new techniques beyond those already used.

% ----------------------------------------------------------
\section*{Acknowledgements}
% ----------------------------------------------------------

The author thanks Omar Khayyam, whose geometric intuition for self-similarity
across scales echoes in the nesting descent at the heart of this proof.

Three earlier attempts informed the present work: a topological proof via
finite preservation~\cite{Wakil2025a}, a nesting descent
argument~\cite{Wakil2025b}, and a spiral descent
formulation~\cite{Wakil2025c}.
Each illuminated a different facet of the problem --- in particular, the
spiral descent attempt identified the critical failure of the compressed
$S$-map and pointed toward the simple $S$-map resolution --- before the
present approach crystallised.

The gap in version~3 was identified through rigorous self-review.
The three-lemma structure replacing Corollary~4.2 emerged from that analysis.

This work was carried out entirely within The ARC Institute of Knowware.

% ----------------------------------------------------------
%  Appendices
% ----------------------------------------------------------

\appendix

\section{Verification of the Mod-16 Transition Table}
\label{app:cascade}

The Forced Cascade (Lemma~\ref{lem:cascade}) rests on the following explicit
computations for $S(n)=(3n+1)/2$.
The final column records the number of $S$-steps required to produce an even
value; ``---'' indicates no exit (only $r=15$ survives under additional
forcing).

\medskip
\begin{center}
\begin{tabular}{@{}clcc@{}}
\toprule
Residue $r$ & $S(16k+r)\bmod 16$ & Steps to even & Notes \\
\midrule
$1$  & $8k+2$, always even              & $1$ & $\equiv 0\pmod{2}$ immediately \\
$5$  & $8k+8$, always even              & $1$ & \\
$9$  & $8k+14$, always even             & $1$ & \\
$13$ & $8k+20\equiv 4\pmod{8}$, even    & $1$ & \\
$3$  & $\in\{5,13\}\pmod{16}$           & $2$ & next step is even \\
$11$ & $\in\{1,9\}\pmod{16}$            & $2$ & next step is even \\
$7$  & $\in\{3,11\}\pmod{16}$           & $3$ & two more steps to even \\
$15$ & $\in\{7,15\}\pmod{16}$           & --- & only survivor (see proof) \\
\bottomrule
\end{tabular}
\end{center}

\medskip
\noindent
For $r=15$: $S(16k+15)=24k+23$.
When $k$ even, $24k\equiv 0\pmod{16}$, so $24k+23\equiv 7\pmod{16}$.
When $k$ odd, $24k\equiv 8\pmod{16}$, so $24k+23\equiv 15\pmod{16}$.
Hence $S(16k+15)\in\{7,15\}\pmod{16}$ as stated.
The case $S(q)\equiv 7\pmod{16}$ is eliminated by the closure of $E$ as
detailed in the proof of Lemma~\ref{lem:cascade}.

\section{Verification of the Oddness Lemma}
\label{app:oddness}

For $p=2s$ (even) and $q=32s+15$, the $S$-orbit satisfies:
\[
  S(q) \;=\; 48s+23 \;\equiv\; 7\pmod{16}.
\]
Writing $m$ for the integer satisfying $S(q)=16m+7$ (so $m=3s+1$):
\[
  S^2(q) \;=\; 24m+11.
\]
In all cases $S^2(q)\equiv 3\pmod{4}$ (regardless of the parity of $m$),
so $S^3(q)$ is odd.
Moreover:
\[
  S^3(q) \;=\; 36m+17 \;\equiv\; 1\pmod{4},
\]
regardless of the parity of $m$.
(Direct check: $m$ even gives $36m+17\equiv 17\equiv 1\pmod{4}$;
$m$ odd gives $36m+17\equiv 36+17=53\equiv 1\pmod{4}$.)

Since $S^3(q)\equiv 1\pmod{4}$, Fact~\ref{fact:parity} gives $S^4(q)$
even, contradicting $q\in E$.

\medskip
\noindent\textit{Numerical verification.}
The value $S^4(32s+15)$ is even for every $s\in\{0,1,\ldots,999\}$, as
verified by computer.
The specific mod-16 residues of $S^3(q)$ vary (attaining values in
$\{1,5,9,13\}\pmod{16}$, all $\equiv 1\pmod{4}$), but the conclusion that
$S^4(q)$ is even is uniform across all cases.

% ----------------------------------------------------------
%  Bibliography
% ----------------------------------------------------------

\bibliographystyle{amsalpha}
\begin{thebibliography}{Roos}

\bibitem[Col37]{Collatz1937}
L.~Collatz,
\textit{On certain sequences of integers defined by mappings},
unpublished manuscript, 1937.

\bibitem[Roos]{Roosendaal}
E.~Roosendaal,
\textit{On the $3x{+}1$ problem},
\url{http://www.ericr.nl/wondrous/},
accessed March 2026.

\bibitem[Tao19]{Tao2019}
T.~Tao,
\textit{Almost all Collatz orbits attain almost bounded values},
Forum of Mathematics, Pi \textbf{7} (2019), e7.
\href{https://doi.org/10.1017/fmp.2019.1}{doi:10.1017/fmp.2019.1}

\bibitem[Wak25a]{Wakil2025a}
K.~Wakil,
\textit{The Collatz conjecture: A topological proof via finite preservation},
preprint, November 2025, The ARC Institute of Knowware.

\bibitem[Wak25b]{Wakil2025b}
K.~Wakil,
\textit{The Collatz conjecture: A complete proof via nesting descent},
preprint, November 2025, The ARC Institute of Knowware.

\bibitem[Wak25c]{Wakil2025c}
K.~Wakil,
\textit{The Collatz conjecture: A proof via spiral descent},
preprint, November 2025, The ARC Institute of Knowware.

\end{thebibliography}

\end{document}
